\chapter{Datasets}
\label{app:datasets}

This appendix describes the datasets used in the project, for training the models and for validating the system: the corpus of the final models, the corpus of the first pipeline, the description of each dataset family, and the reasons behind the choice of the validation recordings. Every source is a public academic release, public sample data from a vendor, or traffic shared openly by the community; no dataset was bought, licensed under a restrictive commercial agreement or obtained under a non disclosure agreement.

\section{Corpus of the final models}

The final models were trained and evaluated per vehicle, with a split by recording: some complete recordings for training, and at least one complete recording kept aside as a holdout. The fingerprints of the recordings guarantee that no recording is on both sides. Table~\ref{tab:app-final-corpus} gives the corpus of the four models in production.

\begin{table}[H]
\centering
\footnotesize
\caption{Corpus of the final models in production.}
\label{tab:app-final-corpus}
\begin{tabularx}{\textwidth}{|P{2.6cm}|Y|r|r|r|}
\hline
\textbf{Model} & \textbf{Source} & \textbf{Training recordings} & \textbf{Training frames} & \textbf{Holdout frames} \\
\hline
Jaguar I-PACE & Community EV logs (EV-CANlogs) & 3 & 233{,}598 & 50{,}361 \\
\hline
Kia EV6 & Community EV logs (EV-CANlogs) & 2 & 448{,}198 & 438{,}170 \\
\hline
Nissan Leaf 2019 & CSS Electronics sample data (CANedge2 logger) & 8 & 1{,}317{,}070 & 491{,}969 \\
\hline
ROAD ambient & ROAD dataset, ambient recordings & 10 & 1{,}437{,}195 & 150{,}000 \\
\hline
\textbf{Total} & & \textbf{23} & \textbf{3{,}436{,}061} & \textbf{1{,}130{,}500} \\
\hline
\end{tabularx}
\normalsize
\end{table}

Six other vehicles did not reach production. The MG ZS EV model was held back after its evaluation, because about 78\% of its holdout frames belong to identifiers absent from its only training recording. The Hyundai Ioniq 5, the Kaggle normal traffic dataset and the Tesla have only one recording each, so their models could not be evaluated on a holdout. No final model was built for the Škoda Enyaq, whose recordings are all kept as protected evaluation files, nor for the older Nissan Leaf family, which we excluded because its recordings mix two logger formats.

\section{Corpus of the first pipeline}

The models of the first pipeline (Section~\ref{sec:impl-state}) were trained on a corpus of \textbf{more than seven million \gls{can} frames} from sixty six sources, summarised in Table~\ref{tab:app-training-corpus}.

\begin{table}[H]
\centering
\caption{Main sources of the corpus of the first pipeline.}
\label{tab:app-training-corpus}
\begin{tabular}{|P{4.2cm}|P{6.8cm}|r|}
\hline
\textbf{Source} & \textbf{Description} & \textbf{Frames} \\
\hline
ROAD ambient~\cite{verma2020road} & Ambient recordings (without attacks) of the ROAD dataset & 1{,}737{,}195 \\
\hline
Car-Hacking (normal rows)~\cite{seo2018gids} & Public benchmark, normal traffic only & 1{,}200{,}000 \\
\hline
Nissan Leaf              & Community EV \gls{can} logs & 1{,}163{,}181 \\
\hline
BMW i3                   & Community EV \gls{can} logs & 1{,}044{,}993 \\
\hline
Kia EV6                  & Community EV \gls{can} logs & 886{,}368 \\
\hline
13 other sources         & Kia Soul, Jaguar I-PACE, MG ZS EV, Tesla and others & about 1{,}032{,}329 \\
\hline
\end{tabular}
\end{table}

For the first pipeline, we tried to train sixteen bundles, one per vehicle or dataset. Fourteen gave a usable model. The two others, BYD Atto 3 and Ford Mustang Mach-E, had too few frames in any driving state to reach the 5{,}000 frames per state that this pipeline required. Today, the model selector of the client lists 20 entries: the general model, 14 vehicle models, each with its status (Table~\ref{tab:model-status}), and 5 \gls{dbc} profiles without a model of their own, including the bench rig, which are scored with the general model.

\section{Dataset descriptions}

Tables~\ref{tab:app-css} to~\ref{tab:app-community} describe the four families of datasets used in the project with the same eight criteria: source, license, vehicle, format, purpose, trustworthiness, strengths and limits.

\begin{table}[H]
\centering
\footnotesize
\caption{CSS Electronics sample data (Kia EV6 and Nissan Leaf 2019).}
\label{tab:app-css}
\begin{tabularx}{\textwidth}{|P{3cm}|Y|}
\hline
Source & Published by CSS Electronics, a manufacturer of CANedge data loggers, as free public sample data showing real \gls{can} and \gls{uds} logging~\cite{css_evdata}. \\
\hline
License and access & Freely downloadable from the public pages of CSS Electronics, for evaluation and demonstration. \\
\hline
Vehicles & Kia EV6 (EV Data Pack) and Nissan Leaf 2019, both real, unmodified production vehicles. \\
\hline
Format & \gls{mf4} binary logs (ASAM MDF version 4), converted into the common frame stream by our \gls{mf4} conversion (Section~\ref{sec:impl-mf4}). \\
\hline
Purpose in this project & The one hour Kia EV6 recording is the dataset of the first validation iteration: root cause analysis, the six fixes and all the measurements before and after. The Nissan Leaf 2019 recordings are the training and holdout data of one of the final models. \\
\hline
Trustworthiness & High: recorded by a professional logger on real vehicles. The Kia EV6 recording was made during an actual drive with active \gls{uds} polling, which gave us an independent fact (the vehicle was being driven) to check the state classification of Fix 2. \\
\hline
Strengths & Real production vehicles; frame counts that can be checked against the \gls{mf4} metadata; a \gls{dbc} profile exists for both vehicles. \\
\hline
Limitations & One recording of the Kia EV6 in the first iteration, and one logger for the Nissan Leaf 2019, so the results do not cover every driving situation. \\
\hline
\end{tabularx}
\end{table}

\begin{table}[H]
\centering
\footnotesize
\caption{ROAD dataset~\cite{verma2020road}.}
\label{tab:app-road}
\begin{tabularx}{\textwidth}{|P{3cm}|Y|}
\hline
Source & Published by its authors (Verma et al.) as an open academic benchmark, the ``Real ORNL Automotive Dynamometer'' \gls{can} intrusion detection dataset. \\
\hline
License and access & Released publicly by the authors for research use. \\
\hline
Vehicle & Traffic of a real vehicle on a dynamometer test bench, with ambient recordings (without attacks) and recordings with labelled attacks. \\
\hline
Format & Text log in the candump style with timestamps: \texttt{(ts) can0 ID\#hexdata}. \\
\hline
Purpose in this project & Only the ambient recordings were used: in the corpus of the first pipeline, and as the training and holdout data of the final ROAD ambient model. \\
\hline
Trustworthiness & High, as an academic benchmark built for the evaluation of intrusion detection, although it was recorded on a test bench and not on the road. \\
\hline
Strengths & Documented provenance; a clear separation between ambient and attack recordings, so the ambient part can be used safely. \\
\hline
Limitations & Test bench traffic, not road traffic. We did not use the labelled attack recordings, so this dataset gives us no measured detection rate; evaluating the final models on them is future work (Chapter~\ref{chap:validation}). \\
\hline
\end{tabularx}
\end{table}

\begin{table}[H]
\centering
\footnotesize
\caption{Car-Hacking dataset~\cite{seo2018gids}.}
\label{tab:app-carhacking}
\begin{tabularx}{\textwidth}{|P{3cm}|Y|}
\hline
Source & Published by the Hacking and Countermeasure Research Lab (HCRL) and used in the GIDS paper as a benchmark for intrusion detection. \\
\hline
License and access & Released publicly by HCRL for research use. \\
\hline
Vehicle & \gls{can} traffic of a real vehicle, with labelled normal and attack traffic. \\
\hline
Format & \gls{csv} with timestamp, identifier, length and payload bytes, plus a label column. \\
\hline
Purpose in this project & Only the normal rows were used, in the corpus of the first pipeline, as a second independent public source of normal traffic. The final models do not use it. \\
\hline
Trustworthiness & High, as one of the most cited public benchmarks of \gls{can} intrusion detection; we only used its normal rows. \\
\hline
Strengths & Well known, which makes its statistics easy to compare with published work. \\
\hline
Limitations & One test vehicle in controlled conditions; its attack labels were not used. \\
\hline
\end{tabularx}
\end{table}

\begin{table}[H]
\centering
\footnotesize
\caption{Community EV \gls{can} logs.}
\label{tab:app-community}
\begin{tabularx}{\textwidth}{|P{3cm}|Y|}
\hline
Source & \gls{can} logs shared publicly by communities of electric vehicle owners and researchers, mainly the EV-CANlogs collection~\cite{evcanlogs}, covering the Nissan Leaf, BMW i3, Kia EV6, Kia Soul, Tesla, Jaguar I-PACE, MG ZS EV, Hyundai Ioniq 5, Škoda Enyaq and other vehicles. \\
\hline
License and access & Shared publicly by their contributors; used for training and evaluation only, and not redistributed with this project. \\
\hline
Vehicle & One real vehicle per source. \\
\hline
Format & Mostly \gls{csv} logs in the SavvyCAN style (timestamps in microseconds) or in the candump style, and some Vector (.asc) logs. \\
\hline
Purpose in this project & Most of the corpus of the first pipeline, and the training and holdout data of the final Jaguar I-PACE and Kia EV6 models. \\
\hline
Trustworthiness & Variable, and less documented than the academic benchmarks. For the final models, the evaluation on a holdout recording of the same source makes the results checkable. \\
\hline
Strengths & Wide coverage of vehicles, which is what makes models per vehicle possible. \\
\hline
Limitations & Informal provenance and uneven documentation; few recordings per vehicle, which is why several vehicles have no holdout. \\
\hline
\end{tabularx}
\end{table}

\section{Choice of the validation recordings}

Training and validation had different priorities. For training, we wanted many independent recordings per vehicle, so that a model does not depend on one session. For validation, we wanted recordings that we could check independently.

In the first iteration, a recording was suitable only if it met four conditions: it came from a real, unmodified vehicle; it was long enough for measurements before and after a fix to be meaningful; its frame count could be checked against the metadata of the file, to rule out a parsing defect being taken for a detection defect; and the vehicle already had a profile and a trained model, so that the path specific to the vehicle was tested, and not just the generic fallback.

The Kia EV6 recording of the CSS Electronics EV Data Pack was the only recording available that met these four conditions at once: an unmodified production vehicle, 97{,}741 frames over about one hour, a frame count that could be checked against the channel group metadata of the \gls{mf4} file (which is how we found the conversion defect of Section~\ref{sec:validation-setup}), and a vehicle with its own profile and trained model, the path that Fix 1 showed was broken. Its active \gls{uds} polling was an extra advantage: it produces a structured pattern of requests and responses, which allowed us to explain the timing share of the alerts (Table~\ref{tab:rootcause}) by its structure instead of treating it as noise.

In the second iteration, the rule was simpler and stricter: each final model is evaluated on complete recordings that were never used to fit it or to set its threshold (Table~\ref{tab:app-final-corpus}).

\section{Dataset used to check the separation of sessions}

To check that no data leaks from one session or dataset into another in the explanation module (Section~\ref{sec:validation-ai}), we used a second, independent recording: a candump log recorded on a Renault Clio during a denial of service experiment, from another vehicle and another capture tool, with identifiers unrelated to the Kia EV6 recording (for example 0xC6 and 0x511). We chose it because it shares no vehicle profile, no logger format and no identifier with the Kia EV6 recording, so any mixing of the two would be obvious.
